As file sizes continue to increase, efficient compression algorithms become more important in order to keep up with the demand for transmitting large amounts of information. Images, a crucial datatype found all across the web, is a prime example of data that requires compression, especially as increasing display resolution pushes the need for higher (and more space-demanding) resolution images \cite{jayasankar_survey_2021}. 

Data compression is the process of representing data in a form that occupies less space than the raw data. Compressed data is found everywhere today; almost all files that users interact with are stored in a compressed format. There are two types of data compression: lossless compression and lossy compression. Lossless compression involves representing the original data such that no information is changed or lost, with common examples being \textit{.zip} and \textit{.tar} files. Lossy compression, as the name suggests, permanently alters the information in order to achieve higher compression. In the case of images, lossy compression aims to discard the information that is least perceptible to humans so that compression can be more efficient while maintaining image quality. This way, the information that is lost via lossy compression is information that the human visual perception system does not need in order to recognize the image.

One of the most widely used image compression standards is the JPEG Compression Standard, which has been in use for over 30 years \cite{wallace_jpeg_1992}. This paper deals with the JPEG Compression Standard, exploring its use of different compression techniques and then tweaking its use of transform coding, a crucial part of the compression process implemented in JPEG. Originally utilizing the Discrete Cosine Transform (DCT), this project dives into other transforms by using different transforms in the JPEG pipeline to investigate their effects on various metrics used to determine the efficiency of an image compression algorithm.